\clearpage\begingroup\SourceLeaf{91}\HistoricalHead{86}{L}{PHYSIK. BRIEFWECHSEL.}
\begin{center}[10.]\end{center}
\begin{center}\textsc{Gauss} an \textsc{Encke}.\quad Göttingen, 21. März 1834.\end{center}
\noindent . . . . . . . . . . . . . . . . . . . . .\par
\indent Gestern und heute bin ich zum ersten male in die Betrachtung der stündlichen
Variationen eingetreten. Meistens ist (während der 44 Stunden) von
10 zu 10 Minuten aufgezeichnet, so jedoch dass jeder Ansatz ein Mittel aus
den 5 nächsten runden Minuten ist; z. B. für \(8^{\mathrm{h}}\,50'\) ist beobachtet bei
\(8^{\mathrm{h}}\,48'\,0''\), \(49'\,0''\), \(50'\,0''\), \(51'\,0''\), \(52'\,0''\) an einer Uhr, die genau die M. Z. zeigt.
Der Verlauf am ersten Tage ist eben der gewöhnliche gewesen, allein in der
vorigen Nacht von \(3^{\mathrm{h}}\) bis \(8^{\mathrm{h}}\) haben grosse Anomalien Statt gefunden, auf deren
Parallelismus mit anderen Beobb. ich sehr neugierig bin. Es haben unser 5
die Wachen getheilt und in diesem Augenblick, 21. März \(11^{\mathrm{h}}\) Abends, ist noch
ein Beobachter an der Arbeit. Reisete Hr. \textsc{Höfer} einen Tag später, so würde
ich ihm sogleich die Abschrift mitschicken, die nun später an H[errn] \textsc{Poggendorff}
eingeschickt werden wird.

Eine \textls[120]{vorläufige} Intensitätsbestimmung (mehreres Zubehör dazu ist erst
vor kurzem fertig geworden) in diesen Tagen hat \(1{,}769\) gegeben, fällt also
zwischen die älteren Bestimmungen.

Der Gebrauch grosser Nadeln hat etwas überaus angenehmes. Die Amplitude
der Schwingung nimmt so langsam ab, dass sie erst nach 2 oder \(2\frac{1}{4}\)
Stunden auf die Hälfte kommt.

Die Scala misst einen (ganzen) Bogen von \(7\frac{1}{2}\) Grad; fängt man also mit
so kleinen Schwingungen an, die noch innerhalb der Scala fallen, so sind
die Schwingungen selbst nach 12 Stunden noch gross genug, um sicher beobachtet
zu werden. Hänge ich eine neue Nadel auf von noch unbekannter
Schwingungsdauer, so ist die Beobachtung von einem halben Dutzend Schwingungen
zureichend, um nach 4 oder 5 Stunden, über die Zahl der verflossenen
Schwingungen nicht ungewiss zu sein. Die gegenwärtig aufgehängte Nadel
braucht \(21''\), sie ist aber noch keine der stärksten, die kräftigste würde circa
\(19''\) gebrauchen.

Eine Hauptschwierigkeit ist sich gute Spiegel zu verschaffen. Ein von
\textsc{Repsold} verfertigter fand sich \textls[120]{ganz unbrauchbar}. Es ist jetzt ein Spiegel
50 mm hoch, 75 breit eingesetzt, den wir durch Zerschneidung und Foliirung
\clearpage\SourceLeaf{92a}\HistoricalHead{87}{R}{GAUSS AN ENCKE.}
\noindent eines \textsc{Troughton}schen Glasdachs gewonnen haben. Dieser ist recht gut, es
ist aber ausserordentlich schwer, \textls[120]{gute} Spiegel von solcher Grösse zu machen.
Ein Künstler in Braunschweig beschäftigt sich jetzt mit dieser Aufgabe. Ein
Paar kleinere, die er als Probe hierher gebracht hat, werde ich morgen versuchen.
Der, wie gesagt ganz misrathene von \textsc{Repsold} hatte nur 40 mm bis
60 mm. Noch kleinere, aber für meinen Zweck \textls[120]{viel} zu kleine, haben wir
früher aus München erhalten, die sehr gut sind, allein die Schwierigkeit
wächst mit der Grösse. . . . . . . . . . . . . . . . . . . . . .
\clearpage\endgroup
